\chapter[Why Theatre Makes Us Cry]{Why Make-Believe on Stage Makes Us Cry}
\section{A Piece of Paper Worth Only a Few Notes}
Pick up an object: a theatre ticket.

It may be thin printed paper or merely a QR code on your phone. Once scanned, its job is done. No refund or exchange; a few hours later it may not even deserve keeping as a souvenir. What did your tens, hundreds, or more yuan buy?

A seat, and two hours' permission to watch people pretend to be other people. The king who falls to a sword will rise for his curtain call ten minutes later without blood on his forehead. The mother who collapses in tears will remove her makeup and hurry for the last train. Nothing is real: story, identities, even walls and forests are canvas and light.

Yet something strange happens: you cry. Tears, heartbeat, sweaty palms, and the long speechless silence afterwards are real. You spend genuine feeling on something that never happened, perhaps more than on actual events. Neighbours quarrelling downstairs may leave you cold, while fictional lovers parting choke you up.

Consider the absurdity of the price. If a stranger said, "Give me two hundred yuan and I will make you cry over something invented", you would think them mad. Yet that is precisely what theatres sell, with buyers every night.

Why do unreal events onstage produce real tears? What are we paying and crying for? This mechanism of knowingly feeling for fiction is among humanity's oldest and strangest entertainments. To understand it, return to an early scene.

\section[Ten Thousand Fall Silent]{Ten Thousand People Fall Silent on a Hillside}
Time: an early spring morning in the fifth century BCE. Place: Athens, the Theatre of Dionysus on the Acropolis's southern slope.

You climb through crowds. Today brings the City Dionysia's main attraction: tragic competition. The city takes a holiday, courts close, and reportedly even prisoners receive temporary leave for theatre. Semicircular stone tiers follow the hill, holding thousands; historians estimate over fourteen thousand at its peak, a sizeable share of the population. Adult male citizens predominate, with women, children, and foreigners too. Seating eligibility has its own account, to which we will return.

The stone is cold; some sit on cloaks. Olive oil and figs scent the air as vendors circulate. Below lies a circular dancing space, the orchestra. Behind it stands a temporary wooden building, backstage and scenery together. Today's story takes place before a palace, so the wooden hut is a palace.

A signal sounds. Within seconds, thousands of voices recede into extraordinary silence. Wind and distant Acropolis flags become audible. A masked chorus of a dozen or more enters, singing its processional song around the orchestra. Then the hut opens and a masked actor steps out, voice rolling uphill to the highest row.

Today brings Oedipus, who kills his father and marries his mother, or Medea, who kills her children in revenge. Everyone knows the ending. Myths are limited; endings cannot change; nobody expects a twist. Yet a farmer beside you, who travelled days to attend, clenches his fists at Medea's accusations. Nobles ahead bow their heads during the chorus's song of fate. At the worst moment, crying gives way to a deeper silence: thousands hold their breath for someone imaginary.

Afterwards, noise returns as people dispute which of three competing playwrights deserves first prize. Judging is public business, jurors selected by lot, the prize an ivy wreath, materially cheap but worth poets' fiercest effort.

Remember this hillside. It is not merely a prototype entertainment venue but a civic instrument, gathering a city around shared imaginary suffering to see what happens. We will examine its workings.

\section{How Can Something Unreal Really Hurt?}
Lay out the conflict before answering.

Oedipus does not exist. Even if he had a historical model, the masked hired Athenian onstage is not him. Medea's cries are rehearsed and her words memorised by other actors too. Every spectator, farmer or magistrate, knows it. Nobody is deceived. Yet tears fall.

Later philosophers named this genuine difficulty the paradox of fiction:

\begin{itemize}
\item You must believe something to feel distress over it. Fire downstairs frightens you because you believe it real.
\item You do not believe the onstage events. Knowing they are fictional keeps you from rushing to rescue Medea's children.
\item Yet your distress is genuine, neither feigned nor needing an audience.
\end{itemize}
At least one of these three claims must be wrong. Which?

Some say absorption temporarily makes spectators forget the fiction. Tears show momentary belief. But if they truly believed, surely they would stop the murderer or call authorities. Across more than two thousand years, audiences have remained seated, seemingly holding two thoughts at once: this is unreal, and this matters.

Others say the emotion itself is theatrical, manufactured to suit the atmosphere. But who watches your melancholy walk home alone, or your sudden tears rereading a speech at midnight?

A third question cuts deeper: even if feeling is real, what is its object? Is grief for a nonexistent person the same as grief for a real person? If tomorrow's news described a mother genuinely in Medea's predicament, would you feel as in the theatre? Probably not. News brings anger, shock, and an urge to act. Theatre brings something quieter, concentrated, even strangely satisfying. You may enjoy crying for Medea; crying for an actual person is different.

This pleasurable sadness is theatre's true oddity. It breaks everyday emotional rules. Before continuing, ask yourself whether theatrical tears seem more real or less.

\section[Expelling Poets, Defending Tragedy]{One Philosopher Expels Poets, Another Defends Tragedy}
Debate begins before the hillside applause ends. Theatre's usefulness and dangers prompted the West's first great dispute over art, and its sides still argue today.

\textbf{Position 1: Theatre is bad, or at least dangerous to a good city.}

Plato, Socrates' pupil and one of Western philosophy's fiercest critics, notoriously proposed excluding tragic poets from his ideal Republic. Often mocked as philosophical ignorance of art, his position deserves its strongest case: the argument is substantial.

It has at least three layers. First, theatre imitates imitation. A bed already copies the Form of Bed; a painted bed or performed carpenter copies the copy, twice removed from truth. Spectators dwell on shadows rather than things. Second, drama nourishes our worst parts. Roughly, Plato argues that real grief calls for rational restraint, whereas tragedy rewards wailing and breast-beating. Applauding others' emotional surrender votes for losing control within ourselves, gradually soaking reason's walls soft. Third, we become what we imitate. Good people playing cowards, murderers, or shrews corrupt themselves through convincing performance. He allows hymns to gods and good people, but expels portrayals of wrongdoing.

Notice the structure: Plato does not call theatre unappealing. \textbf{Its very appeal and effectiveness make it dangerous}. Indifference could never produce this prohibition. He fears precisely the instrument we are investigating.

\textbf{Position 2: Theatre is harmless and essential to civic life.}

This is Aristotle, Plato's own pupil. His Poetics almost reads as a point-by-point answer. Imitation is natural, not degrading: children learn speech and action through it, and humans naturally enjoy it. Tragedy does not weaken reason but cleanses emotion, through the famous catharsis we will examine shortly. Teacher and pupil inspect one theatre; one sees poison, the other medicine.

\textbf{Other voices matter too.}

Not all spectators are philosophers. For the farmer who travelled days, Dionysia is an annual holiday and enjoyable drama needs no further justification; usefulness is a question for the overfed. For Athens, theatre is politics: officially sponsored festivals belong to civic life, reportedly with subsidies enabling poor attendance. Performers occupy subtler ground. Ancient Greek actors were admired professionals, while professional entertainers in ancient China, changyou, had low status and could belong to legally degraded categories. Yet histories preserve their talent for counsel disguised as jokes. In the Records of the Grand Historian's "Biographies of Jesters", You Meng impersonates the late minister Sunshu Ao before King Zhuang of Chu, awakening the king to provide for the minister's descendants. Thorough make-believe achieves something real.

Beyond Athens, civilisations answer theatre's purpose differently.

Ancient India gave it higher standing. The Natyashastra, attributed to Bharata and compiled around the centuries surrounding the beginning of the Common Era, raises drama to a fifth Veda, another sacred teaching accessible to all castes. Sanskrit dramatic theory centres on rasa, "flavour": theatre distils feeling into something savoured, compassion, heroism, erotic love, comedy, and more. Spectators taste purified emotion rather than merely plot. Kalidasa's Shakuntala exemplifies this aesthetic: lovers' misunderstanding, a lost ring, a sage's curse, and final reunion.

China's theatre grew in different soil. Mature opera dates to the Song and Yuan, with roots in older yue, the tradition of music and performance. The Great Preface to the Book of Songs describes emotion becoming speech, then sighs when speech fails, singing when sighs fail, and finally unconscious dancing. Notice the progression: speech $\rightarrow$ sigh $\rightarrow$ song $\rightarrow$ dance. Performance is not language's accessory but its emotional overflow. In the Yuan, Guan Hanqing's Injustice to Dou E takes one vulnerable woman's wrongs to cosmic indignation. We will return to it.

Japanese Noh reaches another extreme. In Fushikaden, written early in the fifteenth century, Zeami repeatedly locates moving performance in restraint rather than fullness. The actor barely moves, the mask stays fixed, and emotion lives in pauses and turns: roughly, concealment lets the flower bloom. A Noh actor may take minutes for one step, while experienced spectators find it overwhelming.

Late-sixteenth-century London's Globe housed different ranks together: ordinary people stood for a penny, gentlemen paid more to sit. Shakespeare's As You Like It voices the famous metaphor that the world is a stage and people actors with entrances and exits. If life itself is role-playing, calling theatre false turns upside down: the stage may be the most honest place.

Do not forget genre. Athenians sharply separated tragedy and comedy: patricide and incest in the morning, Aristophanes mocking rulers and gods in the afternoon. Comedy competed officially too; the city funded contests that ridiculed itself. Classical Sanskrit drama usually ends happily, suffering a passage rather than a destination. Chinese opera also favours reunion, though intriguingly: Dou E returns after death as a wronged ghost to secure justice unavailable among the living. Tragedy or happy ending? Tragicomedy eventually refuses the choice. Shakespeare's late Winter's Tale moves from jealousy destroying a household to a statue coming alive. Three centuries later, Samuel Beckett labels Waiting for Godot a tragicomedy: two people await someone who never comes, leaving spectators unsure whether to laugh or cry, exactly its point.

Five traditions, five answers: danger, medicine, sacred teaching, emotional outlet, life's mirror. Is there a common instrument beneath them?

\section[Thinking Bridge: Performance Materials]{A Thinking Bridge: Four Materials of Performance}
Outsiders ask whether a show is enjoyable; insiders inspect how it works. That inspection is a portable tool. Any performance, Greek tragedy or school show, can be divided into four materials: \textbf{dialogue, action, stage space, audience}.

\textbf{First: Dialogue, what is said.}

Drama largely speaks, but its words differ fundamentally from everyday talk: they are prewritten. Your grandmother saying it is cold responds spontaneously; a character saying so repeats rehearsed words written long before. Dramatic dialogue is premeditation disguised as improvisation. Ask whom words address. Apparently another character, actually the audience, every syllable. Hence unusual devices like soliloquies and asides: characters speak thoughts aloud not through madness but because spectators need access.

\textbf{Second: Action, what is done.}

The theatrical body does not lie, but it speaks dialects. Sadness might make a spoken-drama actor crouch clutching their head, a Beijing-opera actor flick a water sleeve with a cry, or a Noh actor tilt their face slightly. Traditional gestures become \textbf{conventions}, agreed codes. Circling a Chinese-opera stage means travelling a thousand miles; an oar signifies boarding a boat; a whip means riding a horse. Spectators understand through shared codes, not literal resemblance. This will be crucial to explaining genuine feeling for fiction.

\textbf{Third: Stage space, where it happens.}

The stage is no neutral container. Athens' semicircular hillside lets spectators see thousands of fellow citizens as well as the play: you watch drama while the city watches you watching. Later European proscenium stages place audiences in darkness before a framed world, pretending spectators do not exist. Traditional Chinese stages open on three sides, with entrances labelled chujiang, "going out as a general", and ruxiang, "entering as a minister". No wall divides life and drama; audiences applaud and drink tea. Spatial design asks whether spectators are participants or outside voyeurs.

\textbf{Fourth: The audience, who watches.}

The most overlooked material may matter most. Theatre's biggest difference from novels lies below the stage: novels are read alone, plays watched together. Feeling spreads. One laugh lowers others' threshold; collective silence suppresses coughs. Audiences are accomplices, not merely consumers. Without their agreement to pretend, the king instantly becomes a hired masked person again.

At your next play, concert, sports meeting, or opening ceremony, inspect these four materials. The instrument operates everywhere; ordinarily nobody points out its parts.

\section{One Sentence in the Poetics}
Read a little primary evidence. Almost every theatrical discussion encounters one short definition.

In chapter six of the Poetics, Aristotle roughly defines tragedy as imitation of a serious, complete action of suitable length, performed by actors rather than narrated, \textbf{achieving catharsis of pity and fear through those emotions}. The Greek word is katharsis.

The first part is straightforward. Tragedy imitates \textbf{action}, not people. Aristotle deliberately prioritises what someone does and what follows, rather than what sort of person they are. Character serves action. Everyday judgement agrees: ultimately, deeds count.

Performance rather than narration also makes sense: a novel can state anger; theatre needs a person visibly doing it. Being told and being shown engage us differently.

Catharsis is harder. Pity and fear are unwelcome in daily life, so why buy them at a theatre? And what does their purification mean?

The Greek term has medical and religious backgrounds: expelling harmful matter from the body, or ritually removing impurity. Aristotle leaves the term without explanation. The Poetics is a set of lecture notes, its later part, reportedly devoted to comedy, lost. For over two thousand years, readers have disputed the gap.

One reading is purgation: pity and fear accumulate like stagnant water, and tragedy safely lets us cry or fear them out, leaving us cleaner. Theatre becomes emotional drainage. Another is purification through refinement: feelings are not discarded but calibrated. Witnessing extreme suffering teaches proportion, preparing us for life without numbness or collapse. Theatre becomes an emotional gym.

We cannot conclusively establish Aristotle's own choice. Either way, he fundamentally opposes Plato. Plato says theatre \textbf{pours} poison in; Aristotle says it \textbf{drains} poison out. One hillside, one kind of tears, opposite pharmacologies.

Chapter nine adds that poetry is more philosophical and serious than history, roughly because history reports particulars, a person and year, whereas poetry concerns universals, what certain people would plausibly do. Fiction turns around: history tells what happened, drama what could happen. Oedipus may never have existed, but intelligence fleeing a prophecy straight into its fulfilment is a real human structure. Invention can be more universal.

\section[Fictional Events, Real Tears]{Knowing It Is Fiction, Why Cry? Three Explanations}
With materials and evidence, compare three explanations of the paradox. As usual, the phenomenon, tears despite knowledge of fiction, is agreed; its cause is disputed.

\textbf{Explanation 1: Immersion brings temporary belief.}

Common sense says absorption makes us forget the theatre: reason clocks out, feeling clocks in. Early in the nineteenth century, Samuel Taylor Coleridge famously called for a "willing suspension of disbelief": knowing fiction is unreal, we voluntarily set doubt aside.

It captures genuine immersion and its emotional force. But Eastern performance offers a decisive counterexample. Fixed Noh masks and conventional Beijing-opera sleeves continually remind viewers of artifice. Nobody forgets. Immersion theory would predict weak effects, yet experienced audiences are deeply moved by conspicuously unrealistic forms. \textbf{Realism is not necessary for emotional power}. Another explanation is needed.

\textbf{Explanation 2: You know it is pretend but play seriously.}

Twentieth-century philosopher Kendall Walton influentially described appreciating fiction as make-believe, using the same mechanism as children's play. A child riding a stick knows it is a stick, yet within play it is a horse. Stages provide props for the rule that we pretend events real. Your distress for Medea is genuine feeling with an object inside the game, so you remain seated, just as a child does not actually ride to work.

This elegantly revises the paradox's first claim: belief is unnecessary; play suffices. It also explains pleasurable fear, as in rollercoasters or horror, where we know the safety boundary. Its weakness is weight. Is the shudder at Dou E's three fulfilled vows really comparable to playing house? The theory explains tears' mechanism, not their \textbf{gravity}.

\textbf{Explanation 3: Ritual means you are not crying alone.}

Shift from stage to seats. On the Athenian hillside, thousands from one city share a day and imaginary suffering. Studying religion, Emile Durkheim described collective effervescence: gathered people attending to one thing amplify feeling and experience something beyond themselves. Tragic contests grew from Dionysian ritual, retaining choruses, masks, annual festivals, and civic participation. Theatre may never have entirely shed ritual's robes.

Your tears are thus not wholly yours: collective silence and civic solemnity magnify them. This explains why watching a recording alone differs from attending the same play live: the script stays, the ritual changes. It supplies weight missing from play theory, but struggles with private emotion while reading a script. Is that a one-person ritual?

Each explanation illuminates part and leaves darkness elsewhere. Honestly, theatrical tears are probably physiological, playful, and ritual all at once, though nobody can yet specify the mixture.

\section[Further Challenge: Script and Play]{A Deeper Challenge: The Script Is Not the Play}
Now enter a deeper theory from theatre studies through an apparently ordinary distinction: \textbf{the script is not the play}.

A script is words on paper. A play is a unique event among bodies in a particular place and evening. Recipe and dish offer an analogy: recipes endure and circulate; meals are consumed and disappear. Shakespeare's scripts remain in print, but none of his performances survives. Tonight's Hamlet and the Globe's four centuries ago share a text but differ almost entirely. Researchers distinguish \textbf{script}, the written literary text, from \textbf{performance text}, everything actually occurring, words, actions, lights, coughs, rain that evening. One is a recipe, the other a once-only encounter.

This explains why a star's Dou E electrifies while an amateur's identical lines induce sleep: different performance texts. Zeami's restrained flower lives in pauses and turning speeds unreadable on paper, precisely performance's core assets. Where is the play? The script is in a library; the play exists only that evening.

The theory's sharpest edge comes from someone deliberately reversing it. German playwright Bertolt Brecht objected to traditional theatre's emotional routine: absorb, empathise, cry, return home satisfied. Aristotelian cleansing, he argued, anaesthetises reality. Cry for Dou E safely and the world seems tolerable afterwards; anger has drained away. He therefore creates a deliberate \textbf{estrangement effect}, Verfremdungseffekt. Actors address spectators, machinery remains visible in bright light, songs interrupt plots, titles disclose endings. Keep audiences alert rather than crying: do not merely pity people, ask what made their condition.

Our question reaches its depth: is immersion an ability or an anaesthetic? Aristotle says tears cleanse; Brecht says they disarm. This is not ancient history. Every film you label tear-jerking and then forget stands at the dispute's centre.

Avoid another trap: estranging theatre is not automatically superior, nor moving theatre necessarily anaesthetic. Brecht's own classics, such as Life of Galileo, deeply move audiences. His best work activates thought and feeling together rather than eliminating emotion. Estrangement asks tears to pass through the mind. Great theatre often pulls us repeatedly into and out of immersion.

\section{Today: A Theatre in Your Pocket}
This ancient question now performs in your pocket at unprecedented scale.

You may rarely visit theatres but watch performances daily: scripted short videos, twist-filled mini-dramas, livestreamed moods, esports commentary. Personas, plots, acting, sometimes scripts: all theatre. Athenians watched a few festival days yearly; you scroll for hours daily. Dionysia now updates every day.

Our four materials clarify the change.

Dialogue and action: why are mini-drama performances exaggerated? Tiny screens and fractured attention compress conventions to their limit, conflict every three seconds, reversals every ten. It is not simply poor acting but changed stage space. Hillside becomes palm; ritual folds into algorithms. Beijing-opera conventions took centuries; short-video conventions change generations within two years.

The audience changes most intriguingly. Scrolling comments bring the Athenian hillside to a phone. Alone, you still see thousands of simultaneous reactions: technology simulates collective effervescence. Concerts and matches do it directly. Ten thousand voices singing one song genuinely experience Durkheim's beyond-the-individual feeling, with an idol replacing Dionysus. Official fan chants, fixed words at fixed moments, resemble a chorus. Festival processions and concert responses are versions of one instrument.

Aristotle versus Brecht becomes a platform question. Algorithms know your accumulated grievances and supply satisfying, tearful, three-minute revenge. Cry, swipe, next. Afterwards are you clearer-minded or merely more compliant? Plato returns in new clothes: what happens to reason's walls when citizens spend three hours daily absorbed in copies of copies? Yet do not award him victory too quickly. The same screen brings Teahouse to small-town teenagers who never entered a theatre and reunites regional opera with scattered audiences. The instrument itself is neither good nor bad; the issue is who turns its knobs.

A new variable: there may be nobody onstage. Virtual streamers and AI characters lack actors removing makeup and catching trains, perhaps any living performer at all. Play theory finds no problem: if a stick can be a horse, code can be a character. Ritual theory finds something uncanny: ten thousand participants gather around a nonexistent figure on the altar. Whom do tears for that figure address?

\section{Your Turn: What Are Tears Worth?}
Return to the ticket.

With four materials, three explanations, and twenty-four centuries of debate, answer without fence-sitting, but with reasons:

\textbf{Which tears are worth more: those shed for fictional characters in a theatre, or those shed for real suffering in life?}

Before deciding, put these weights on the scales.

Fiction has a strong case. Aristotle calls poetry more universal than history: one real sufferer has a particular fate, while Dou E concentrates everyone whose appeals go unanswered. Your tears may stand for millions. Sanskrit theory distils emotion into rasa, compassion without self-interest, panic, or practical calculation. And the hillside makes private crying civic rehearsal for bearing suffering together.

Reality's case is equally strong. Plato's concern persists: comfortable theatre tears can leave us feeling already virtuous, spending sympathy onstage and withholding it outside. Brecht asks whether discharged emotion is displaced action. Most simply, Medea does not need you; the freezing person on the corner does. How should we reckon fully paid fictional compassion alongside real debts unpaid?

Go deeper and the issue grows harder. If tears' quality depends not on their object's reality but on what they change in you, compare a play making you more patient with others and genuine news provoking three days' anger followed by nothing.

There is no standard answer. But weighing it already makes you more than an ordinary ticket buyer. From now on, whether theatre, cinema, or phone brings tears, you know another layer: how the instrument works and which channel carries your feeling. Knowledge need not diminish tears, but can clarify why you cry.
